\section{K2 是乔戈里峰：基础模型、效率与 Agentic 能力}

本章进入 K2。杨植麟把 K2 的目标拆成几个部分：第一，希望它是非常好的基础模型；第二，希望最大化使用每一份数据，也就是提高 token efficiency；第三，希望它有好的 Agentic 能力，能从“缸中之脑”走向多轮工具使用；第四，开源让能力进入生态。

\begin{figure}[H]
\centering
\includegraphics[width=0.96\textwidth]{k2-design-goals.png}
\caption{Kimi K2 设计目标：基础模型、token efficiency、Agentic 能力和开源同时成立。自制概念图，依据 00:24:58--00:54:08 对谈内容整理。}
\end{figure}

\begin{knowledgebox}{读图：K2 不是单点优化}
基础模型决定底座，token efficiency 决定同样数据能长出多少能力，Muon 和 Rephrase 是训练效率手段，Agentic 能力决定模型能否使用工具，开源决定生态扩散。
\end{knowledgebox}

\subsection{Token efficiency：同样数据，脑子长得更多}

数据墙之后，token efficiency 变得更重要。杨植麟的说法是，喂一样多的数据，希望“脑子”长得更多。K2 会对数据做 Rephrase，也关注 Muon 优化器等方向。这些技术不只是为了 benchmark，而是为了在高质量数据越来越稀缺时，提高每个 token 的训练收益。

\begin{figure}[H]
\centering
\includegraphics[width=0.96\textwidth]{token-efficiency-loop.png}
\caption{Token Efficiency 闭环：数据改写、优化器和训练策略提高每个 token 的收益。自制概念图，依据 00:27:28--00:32:44 对谈内容整理。}
\end{figure}

\begin{knowledgebox}{术语消化：训练效率相关概念}
\begin{tabularx}{\textwidth}{p{0.22\textwidth}p{0.32\textwidth}X}
\toprule
概念 & 作用 & 风险 \\
\midrule
Token efficiency & 同样 token 带来更多能力 & 难以单独度量，常被其他变量混淆。 \\
Rephrase & 改写数据以增加表达变化 & 改写质量差会引入噪声。 \\
Muon & 优化器方向，可能提升训练效率 & 训练可能不稳定，访谈中提到“会炸”。 \\
Data mixture & 数据配比和课程 & 配错会伤害能力或泛化。 \\
\bottomrule
\end{tabularx}
\end{knowledgebox}

\subsection{Agentic 模型的泛化挑战}

上一节讲 K2 的训练效率，本节转向 Agentic 能力的核心难题：泛化。对于 Agentic 模型，最大挑战是泛化。训练环境里的工具、任务和反馈如果太窄，模型可能只会在熟悉场景里表现好；真正有价值的是遇到新工具、新任务、新环境时仍能多轮探索、调用工具并修正行为。

\begin{figure}[H]
\centering
\includegraphics[width=0.94\textwidth]{agentic-generalization.png}
\caption{Agentic 模型的泛化挑战：会用工具不等于能泛化到新任务。自制概念图，依据 00:32:50--00:38:00 对谈内容整理。}
\end{figure}

\begin{importantbox}{Agentic 泛化的判断}
一个模型是否真正 Agentic，不是看它是否会调用一个固定工具，而是看它能否在未知任务中理解工具说明、规划多步动作、从反馈中修正，并把能力迁移到新环境。
\end{importantbox}

\begin{knowledgebox}{Agentic 泛化的三层测试}
\begin{tabularx}{\textwidth}{p{0.22\textwidth}p{0.34\textwidth}X}
\toprule
层级 & 测什么 & 失败样式 \\
\midrule
工具泛化 & 新 API、新命令、新文档是否能理解 & 只会训练中见过的工具。 \\
任务泛化 & 新目标和长程依赖是否能分解 & 计划漂亮但无法落地。 \\
反馈泛化 & 错误返回后能否自我修正 & 一错到底，或重复同一动作。 \\
\bottomrule
\end{tabularx}
\end{knowledgebox}

\subsection{本章小结}

K2 的核心是多目标平衡：强基础模型、训练效率、Agentic 泛化和开源生态必须同时成立。任何一个环节短板，都会限制它从模型变成系统。

